\section{Completion of the conic-image leaf}
\label{app:quartic-conic-completion}

This appendix treats the three conic representatives not covered by the
invariant-field proof in the main text.  The orbit with fixed factor
\(G=z^2\) is handled by the separate exact checker distributed with this
packet.  Here we close the remaining two representatives
\[
 H_4=x^2(x^2,xy,y^2),\qquad
 H_4=xy(x^2,xy,y^2).
\]
The computations retain arbitrary quadratic layer and arbitrary linear part.
They use exact rational arithmetic, no random specialization, and no division
by a parameter before entering the corresponding nonzero chart.

Put
\[
 n=(y^2,-2xy,x^2)^T,
 \qquad
 \det(JH_1+\epsilon JH_2+\epsilon^2JH_3+\epsilon^3JH_4)
 =\sum_{j=0}^9D_j\epsilon^j.
\]
For either fixed factor, the first normal equation gives
\(M=n\cdot H_3\in k[x,y]_5\).  Write
\[
 H_3=T(A,B)+N(M),\qquad
 T(A,B)=(2xA,yA+xB,2yB),
\]
where \(A,B\) are quadrics, and for
\(M=\sum_{i=0}^5m_ix^{5-i}y^i\) take
\[
 N(M)=
 \left(m_4xy^2+m_5y^3,0,
 m_0x^3+m_1x^2y+m_2xy^2+m_3y^3\right).
\]

\begin{theorem}[The double-line conic orbit]
\label{thm:quartic-conic-x2}
No quartic Keller counterexample has
\[
 H_4=x^2(x^2,xy,y^2).
\]
\end{theorem}

\begin{proof}[Exact coefficient proof]
Write
\[
 A_z=a_2x+a_4y+a_5z,
 \qquad
 B_z=b_2x+b_4y+b_5z.
\]
The compatibility equations of \(D_7=0\) first give
\[
 a_4=a_5=b_5=0,
\]
and
\begin{align}
4a_3(a_2-b_4)+(4b_4-a_2)m_4+6b_2m_5&=0,
\label{eq:x2-routing-one}\\
m_5(a_2-2b_4)&=0.
\label{eq:x2-routing-two}
\end{align}
The coefficients
\begin{align*}
[x^4z^2]D_6&=24b_2^2b_4,\\
[x^3yz^2]D_6&=-48b_2b_4(a_2-b_4),\\
[x^2y^2z^2]D_6&=24b_4(a_2-b_4)^2,\\
[y^5z]D_6&=-6b_4m_5(a_2+b_4)
\end{align*}
together with \eqref{eq:x2-routing-one}--
\eqref{eq:x2-routing-two} give four exhaustive stabilizer types:
scalar, rank-one semisimple, nilpotent, and zero.

In the scalar chart normalize
\[
 A=a_0x^2+a_1xy+xz,
 \qquad
 B=b_0x^2+b_1xy+b_3y^2+yz.
\]
The equations give \(M=0\), \(a_1=b_3\), and then
\[
 c_6=b_0(a_0+b_1),\quad
 c_0=a_0^2-2b_0b_3-b_1^2+c_7,
 \quad c_1=4b_3(a_0-b_1).
\]
After the preceding linear solves put
\begin{align*}
 A_*&=4a_0b_0b_1+2b_0^2b_3-4b_0b_1^2-2b_0c_7+\ell_7,\\
 B_*&=2a_0^2b_1-4a_0b_1^2-a_0c_7+2b_1^3+b_1c_7+\ell_4,\\
 S&=2a_0b_0^2b_3+a_0\ell_7-2b_0^2b_1b_3-2b_0\ell_4-b_1\ell_7.
\end{align*}
The exact residual identities are
\[
 D_2=A_*^2x^2-4A_*B_*xy+4B_*^2y^2,
 \qquad
 \det L=S^2,
\]
and
\[
 S=(a_0-b_1)A_*-2b_0B_*.
\]
Thus \(D_2=0\) implies \(\det L=0\).

In the rank-one semisimple chart normalize
\[
 A=a_0x^2+a_1xy+xz,
 \qquad
 B=b_0x^2+b_1xy+b_3y^2.
\]
The first lower equations give
\[
 m_0=0,
 \quad c_{16}=\frac38m_1,
 \quad c_{10}=b_1+\frac5{16}m_2,
 \quad b_3=-\frac9{16}m_3.
\]
Four coefficients of \(D_5\) force \(m_1=m_2=m_3=0\).  The remaining
\(D_4\)-equations are
\[
\begin{aligned}
X&=a_1b_0^2-a_1c_{12}+2b_0b_1^2-b_1c_{13}+\ell_7=0,\\
Y&=a_0a_1b_0+a_0b_1^2+a_1b_0b_1-a_1c_6-b_1c_7+\ell_4=0.
\end{aligned}
\]
Substitution gives \(\det L=0\).

In the nilpotent chart a source translation gives
\[
 A=a_0x^2+a_1xy,
 \qquad
 B=xz+b_3y^2.
\]
The equations reduce to
\[
\begin{gathered}
 m_3=-2b_3,
 \quad c_{15}=m_2^2/4,
 \quad c_9=(a_1+b_3)m_2/2,\\
 b_3(a_1+b_3)=0,
 \qquad
 (a_1+b_3)(3a_1m_2+2a_0b_3)=0.
\end{gathered}
\]
If \(a_1+b_3=0\), the remaining linear equations give \(\det L=0\).
Otherwise \(b_3=m_2=0\); normalize \(a_1=1\).  Then \(D_4=0\) forces
\[
 m_0=m_1=0,
 \quad \ell_6=a_0c_{12},
 \quad \ell_3=a_0c_6,
 \quad \ell_0=a_0(c_0-a_0^2),
\]
and again \(\det L=0\).

In the zero chart \(H_3\) is binary and
\[
 \partial_zH_2=(2px,py+qx,2qy).
\]
The two nonzero stabilizer orbits are represented by \((p,q)=(1,0)\) and
\((0,1)\).  They respectively give the fixed coefficients
\[
 [x^2y^2z]D_5=-8,
 \qquad
 [x^4z]D_5=-8.
\]
Thus \(p=q=0\), so every nonlinear layer is binary.  A target linear change
reduces the map to
\[
 (f_1(x,y),f_2(x,y),z+f_3(x,y)),
\]
whose plane Keller pair has degree at most four and is automorphic by the
low-degree plane theorem used elsewhere in Program~2.
\end{proof}

\begin{theorem}[The nodal two-line conic orbit]
\label{thm:quartic-conic-xy}
No quartic Keller counterexample has
\[
 H_4=xy(x^2,xy,y^2).
\]
\end{theorem}

\begin{proof}[Exact coefficient proof]
The compatibility equations of \(D_7=0\) give
\[
 a_4=a_5=b_2=b_5=0,
 \qquad
 m_0(3a_2-b_4)=m_5(a_2-3b_4)=0.
\]
The pure coefficients
\[
\begin{aligned}
[x^5z]D_6&=-3m_0(3a_2^2+b_4^2),\\
[y^5z]D_6&=-3m_5(a_2^2+3b_4^2),\\
[x^2y^2z^2]D_6&=12(a_2-b_4)^2(a_2+b_4)
\end{aligned}
\]
show that a nonzero \(z\)-matrix has exactly two stabilizer types,
\((a_2,b_4)=(1,1)\) and \((1,-1)\).

In the scalar chart, the next equations force
\[
 a_3=b_0=m_1=m_2=m_3=m_4=0.
\]
After the linear solves, \(D_4=0\) gives
\[
\begin{aligned}
 c_{16}&=4(b_1-a_0),& c_{10}&=2(a_1-b_3),\\
 c_6&=-2a_0(a_0-b_1),&c_9&=2b_3(a_1-b_3),\\
 c_{15}&=-a_0a_1-a_0b_3-a_1b_1+3b_1b_3+c_7.
\end{aligned}
\]
Set
\begin{align*}
 A_*={}&2a_0^2a_1-4a_0^2b_3+4a_0b_1b_3-2a_0c_7
       -2a_1b_1^2+2b_1c_7-\ell_7,\\
 B_*={}&a_0a_1^2-2a_0a_1b_3+a_0b_3^2+a_1^2b_1
       -2a_1b_1b_3-a_1c_7+b_1b_3^2+b_3c_7+\ell_4,\\
 S={}&2a_0^2a_1b_3-2a_0^2b_3^2-4a_0a_1b_1b_3
       +4a_0b_1b_3^2+2a_0\ell_4\\
 &\qquad{}+2a_1b_1^2b_3+a_1\ell_7-2b_1^2b_3^2
       -2b_1\ell_4-b_3\ell_7.
\end{align*}
Then
\[
 D_2=A_*^2x^2+4A_*B_*xy+4B_*^2y^2,
 \qquad \det L=S^2,
\]
and
\[
 S=-(a_1-b_3)A_*+2(a_0-b_1)B_*.
\]
Thus \(D_2=0\) gives \(\det L=0\).

In the anti-scalar chart, the high equations give
\(m_1=4b_0,m_4=4a_3\) and then \(b_0=a_3=0\).  The residual coefficient is
\[
 [xyz^3]D_5=-64,
\]
so the chart is empty.

In the zero chart, again
\[
 \partial_zH_2=(2px,py+qx,2qy).
\]
The stabilizer of \(xy\) has nonzero vector orbits represented by
\((1,0)\) and \((1,1)\).  The first has
\([xy^3z]D_5=-8\).  The second has
\[
 [x^3yz]D_5=-8,
 \quad [x^2y^2z]D_5=16,
 \quad [xy^3z]D_5=-8.
\]
Thus \(p=q=0\), and the all-binary plane reduction proves automorphy.
\end{proof}

\begin{corollary}[All conic-image leaves]
\label{cor:quartic-all-conic-orbits}
Subject to the same low-degree plane input used in the reader manuscript,
no quartic Keller counterexample has a nondegenerate conic as the
projective image of its leading homogeneous part.
\end{corollary}

\begin{proof}
The four reduced representatives in the main conic theorem, together with
\cref{thm:quartic-conic-x2,thm:quartic-conic-xy} and the separate
\(G=z^2\) exact proposition, exhaust the seven parabolic conic orbits.
\end{proof}

\begin{remark}[Evidence boundary]
The supplied SymPy programs reconstruct the full determinant arc with all
lower coefficients present.  They are exact, assertion-based calculations,
but they constitute one computer-algebra lineage rather than an independent
second-system reproduction.  This corollary closes the conic leaf; it does
not establish the exhaustiveness of the global quartic case tree.  The proper
rational-cubic leading-image leaf remains separate.
\end{remark}
